% Source inspection records: research/2026-09-22-pass11-prior-art.md and
% research/2026-09-22-pass15-assessment.md. Glazebrook 1992/1993 were inspected
% at abstract/metadata level only; no theorem-level separation is asserted.
\section{Related work and interpretation}
\label{sec:related}

\subsection{Deteriorating work and initial preparation}

Scheduling with deterioration already asks when delaying or interrupting a
job increases the work required to complete it. Glazebrook
\cite{Glazebrook1992,Glazebrook1993} studies sufficient conditions for
nonpreemptive or permutation policies in stochastic deterioration models.
A more directly inspectable repair-state formulation is given by Gehlot,
Sundaram, and Ukkusuri \cite{GehlotSundaramUkkusuri2020}. In their discrete
model, the selected component gains a fixed repair increment, other components
lose fixed health increments, and both failure and full repair are absorbing.
Their Theorem~1 removes preemptions when each deterioration increment is at
least the corresponding repair increment. Its proof postpones a partially
repaired component and compares the resulting service durations. Their
subsequent model incorporates precedence constraints and a finite deadline
\cite{GehlotSundaramUkkusuri2020Interdependent}.

Our supporting serial reduction uses the same broad idea that service moved
later suffers less deterioration. Its hypotheses differ: zero preparation is
recoverable, controls are continuous and divisible, loss depends on current
preparation, and completed modules may increase available capacity. We do not
present nonpreemption under deterioration as a new general principle. The
main optimization here concerns the \emph{initial preparation vector} and its
ongoing maintenance cost, after an exact exit-time characterization has been
established.

Controllable-processing-time scheduling is consequently another essential
comparison. Shioura, Shakhlevich, and Strusevich
\cite{ShiouraShakhlevichStrusevich2015} minimize linear compression cost under
a common deadline and develop algorithms for linear optimization over a
box-constrained submodular region. Their Section~3.1, Theorem~2, states the
standard greedy rule that contains our common-decay fixed-order optimization.
Indeed, if an order has deadline inequality
\[
 A_\pi-\sum_r w_r p_{\pi_r}\leq X,
 \qquad \sum_r w_rM_{\pi_r}=A_\pi-1,
\]
then the substitution $y_r=w_r(M_{\pi_r}-p_{\pi_r})$ gives
\[
 \max\sum_r\frac{\gamma}{w_r}y_r,
 \qquad 0\leq y_r\leq w_rM_{\pi_r},
 \qquad \sum_r y_r\leq X-1.
\]
Thus one-partial-coordinate existence in this specialization is ordinary
linear allocation. The strict ordering of the model's weights supplies the
additional preparation-prefix property. Taking a finite minimum over orders
does not change the attribution of the fixed-order optimization.

Wei, Wang, and Ji \cite{WeiWangJi2012} combine start-time deterioration with
controllable processing times of the form $a_j+bt-\theta u_j$. Their
Section~3 expands a prescribed order's completion time into a weighted affine
function of resource allocations. For identical module sizes $M$, common
decay $\gamma$, and zero releases, our serial recurrence becomes their
recurrence after setting $D=s-\gamma M>0$, $z=e^{\gamma t}-1$,
$b=\gamma M/D$, $a_j=\gamma M/D$, $\theta=\gamma/D$, and $u_j=p_j$.
This is an exact restricted timing correspondence; their formulation assumes
nonpreemption and optimizes weighted scheduling and resource costs.

Nonlinear resource dependence alone also does not distinguish the present
problem. Li, Wang, and Wang \cite{LiWangWang2015} consider durations
$((a_j/u_j)^k+bt)r^\alpha$ and derive optimal investment and sequencing from
separable convex resource terms and rearrangement. Our initial preparation
instead interacts with waiting through
\[
 \tau_i(t,p;b)=\frac{1}{\gamma_i}
 \log\frac{b-\gamma_i p e^{-\gamma_i t}}{b-\gamma_iM_i}.
\]
The mixed derivative with respect to $t$ and $p$ is strictly positive, whereas
their duration is additively separable in start time and investment. Directly
identifying their resource variable with initial preparation therefore does
not give this model.

A directly inspectable fixed-deadline nonlinear allocation appears in
Choi and Park \cite{ChoiPark2023}, Section~3, Lemma~3. Its proof minimizes
$\sum_j c_j u_j$ subject to $\sum_j(w_j/u_j)^k\leq\delta$, with positive
variables and coefficients. Convexity and the KKT conditions give
$u_j=(\lambda k w_j^k/c_j)^{1/(k+1)}$; the investments can all be interior.
Our fixed-order deadline geometry need not be convex. For example, with
$s=3$, zero releases, $M_1=M_2=1$, $(\gamma_1,\gamma_2)=(2,1)$ and
$H=\log2$, order $1,2$ requires
\[
 y\geq3\sqrt{3-2x}-4.
\]
The boundary points $(13/25,1/5)$ and $(3/8,1/2)$ are feasible, while their
midpoint $(179/400,7/20)$ is not: squaring the required positive-side
inequality would give $7578\leq7569$. An affine identification with that
convex investment subproblem is therefore impossible. This comparison
excludes that specified reduction; it does not exclude arbitrary coupled
transformations or establish priority over all investment scheduling results.

The unequal-rate concentration theorem is the structural question beyond
these fixed-order correspondences. Its deadline-preserving pair exchange
excludes two interior coordinates even without a common affine clock. The
cost coefficients are the decay coefficients themselves, rather than arbitrary
investment prices. Nonnegative capacity release preserves the exchange; the
result also includes zero releases. The quadratic-loss example identifies a
boundary of this particular concentration property, while leaving seriality
intact. It is not a claim that nonlinear resource models generally concentrate
investment.

\subsection{Maintenance, service, and initialization}

The passage from minimum maintenance to long-run optional throughput is a
storage argument. With $Q=\sum_i p_i$, a ready trajectory under maximum loss
satisfies $\dot Q\leq s-u-C(H)$. Integrating and using bounded storage gives
the average-throughput bound. This is the established connection between a
storage inequality and optimal steady operation discussed by Faulwasser
et al.\ \cite{FaulwasserEtAl2017}. We prove the particular reflected-system
inequality directly; the contribution under study is the readiness geometry
and its optimizer, not the averaging principle.

The model assumes an independent receiver and an instantaneous transfer
primitive. This is stronger than possession of a copied image. Practical
fault-tolerant execution must preserve externally acknowledged obligations:
Scales, Nelson, and Venkitachalam
\cite{ScalesNelsonVenkitachalam2010}, Sections~2.2--2.3, gate external output
on sufficient backup replay information and treat takeover separately.
Live migration likewise distinguishes pre-copy, source suspension, remaining
transfer, commitment, and activation \cite{ClarkEtAl2005}, Section~3.3.
These mechanisms demonstrate why bandwidth accounting alone does not imply
correct ownership transfer or a bounded service response.

Here proportional loss is a pathwise uncertainty envelope. Randomly located
record updates may motivate an expected proportional drift, but that
expectation does not establish the envelope for every update history. A
finite-delay implementation would also need an explicit contract for
acknowledgements, in-flight information, queuing, and fencing. The results
therefore describe perishable preparation under the stated ideal handoff
interface, rather than an implementation-independent migration bound.

Finally, maintaining an optimal state and reaching it from a cold system are
different questions. Paid initialization makes the stationary throughput
construction feasible whenever its maintenance fits the budget. Without that
allowance, strict maintenance slack permits finite warmup, while equality
requires a separate reachability analysis. The critical three-module
construction further separates finite entry into an indefinitely ready
trajectory from finite arrival at a stationary optimum: the former can be
possible when the unique optimum is unreachable. No readiness guarantee is
asserted during an unprepared warmup. This distinction is part of the service
contract, not an adjustment to the long-run objective.

\subsection{Why the charged resource and normal-mode restriction matter}

The source-only upkeep bound need not survive additional independent input.
For one module with $M=\gamma=1$, $s=2$, $a=0$ and $H=\log(3/2)$,
scalar comparison gives $F(p)=\log(2-p)$ and hence $C(H)=1/2$.
The declared model has optional optimum $3/2$. If an independently paid input
of rate $1/2$ can maintain the same preparation, initialize at $p=1/2$ and
use that input while assigning the entire source budget to optional work.
Every request still meets $H$ under the original rate-two exit strategy, and
normal optional throughput is two. Total maintenance has not vanished; its
cost has moved outside the source objective.

Likewise, if permanent transfers are allowed before the request, preserve
the same required service, and have no separately charged objective cost,
the initialized problem admits an immediate counterpolicy: initialize every
module fully and transfer them all. Every later request then has zero exit
time and needs no source preparation. A positive source-upkeep lower bound
therefore requires the stated normal-mode restriction or an objective that
accounts for the admissible alternatives. Neither restriction follows merely
from calling the eventual receiver independent. These counterpolicies delimit
the ideal service class; they do not assert that every physical receiver can
supply either additional control.
